\usepackage[ngerman]{babel}

% Context sensitive quotation facilities
\usepackage[autostyle,german=quotes]{csquotes}

\usepackage{mathptmx}
%\usepackage[T1]{fontenc}


% Erlaubt Positionierung von Abbildungen
\usepackage{float}
% mit dem command \fbox vor dem \includegraphics{} wird das Bild schwarz umrandet



% Erlaube Zugriff auf 'author', 'date', 'title' durch \@author etc.
% Siehe auch:
% https://tex.stackexchange.com/questions/10130/use-the-values-of-title-author-and-date-on-a-custom-title-page
\makeatletter

% fontspec wird statt fontenc verwendet (fontenc ist veraltet)
% fontspec erlaubt Zeichen wie das "ß", ändert aber die Schriftart
% https://tex.stackexchange.com/questions/28450/xetex-%C3%9F-compiles-as-ss-german-umlauts-work
\usepackage{fontspec}
% Überschreibe den Default font
\setmainfont{Times New Roman}

% Entferne Farben, die sonst in Hyperlinks/Referenzen auftauchen würden
\usepackage[hidelinks]{hyperref}

% Erstellt Bookmarks in PDFs und erstellt Hyperlinks von Links
\usepackage{bookmark}

% Global - Unterdrücke Einrückungen bei neuen Absätzen
\setlength\parindent{0pt}
% mit \par wird ein neuer Absatz eingerichtet, welche sofern oben geschriebenes nicht aktiviert ist eingerückt wird. durch  \\ kann ein Satz direkt
% direkt in die neue Zeile gesetzt werden ohne Einrücken, auch wenn automatischen Einrücken nicht unterdrückt ist!

% Enhanced support for graphics
% Im img Ordner befinden sich alle Bilder
\usepackage{graphicx}
\graphicspath{ {./img/} }

% Sorgt dafür, dass Abbildungsunterschriften Linksbündig sind
\usepackage{ragged2e}
\usepackage[format=plain,
      justification=RaggedRight,
      singlelinecheck=false]
     {caption}

% Erlaubt clevere Referenzierung
% http://tug.ctan.org/macros/latex/contrib/cleveref/cleveref.pdf
\usepackage{cleveref}

% Erlaubt die Verwendung von \titleformat
% (LaTeX fügt standardmäßig 'Kapitel 1' etc. vor den eigentlichen Kapitel)
% (Deshalb müssen wir das Titleformat manuell ändern)
% Siehe auch:
% http://www.ctex.org/documents/packages/layout/titlesec.pdf
% \chapter{Überschrift} > Kapitel 1
% \section{Unterüberschrift} > Kapitel 1.1
% \subsection{Unterunterüberschrift} > Kapitel 1.1.1
% \subsubsection{Unterunterunterüberschrift} > Kapitel 1.1.1.1
\usepackage{titlesec}
\titleformat{\chapter}{\normalfont\bfseries\Huge}{\thechapter.}{10pt}{}

% Passt Abstände der Überschriften an
% \titlespacing{command}{left spacing}{before spacing}{after spacing}[right]
\titlespacing*{\chapter}{0pt}{-40pt}{15pt}
\titlespacing*{\section}{0pt}{0pt}{10pt}

% Ermöglicht die Einstellung von Seitenrändern und ähnlichem 
% https://www.namsu.de/Extra/pakete/Geometry.html

\usepackage[a4paper,left=2cm,right=2.5cm,top=2cm,bottom=2.5cm]{geometry}

% Änderung der Schriftgröße :Changing the font size locally
% \Huge \huge
% \LARGE \Large \large
% \normalsize (default)
% \small
% \footnotesize
% \scriptsize
% \tiny
%
% damit nicht der ganze folgende Text, sondern nur bestimmter Text in der Größe geändert wird:
% {\huge Hier den Text einfügen}#


% zum Anzeigen des Seitenlayouts 
%\usepackage{showframe}

% Literaturverzeichnis
% Siehe auch:
% https://www.heise.de/tipps-tricks/LaTeX-Literaturverzeichnis-erstellen-so-klappt-s-4401420.html
\usepackage[backend=biber,bibstyle=authortitle,citestyle=authoryear]{biblatex}
\addbibresource{lit.bib}
\addbibresource{lit-web.bib}


% Multipliziert den Zeilenabstand
\linespread{1,5}

% So macht man direkte Zitate:
% 
%\begin{quotation}
%    \glqq Arguing that you don't care about the right to privacy because you 
%    have nothing to hide is no different than saying you don't care about free 
%    speech because you have nothing to say.\grqq\footnote{\cite{snowden2015}}
%\end{quotation} 

% Ermöglicht eine Veränderung der Zeilenabstände bei einer Nummerierten Aufzählung
% \setlist[enumerate]{itemsep=xcm}
% \begin{enumerate}
\usepackage{enumitem}

% subsubsection wird jetzt nummeriert und im Inhaltsverzeichnis angezeigt
% \setcounter{tocdepth}{3}
% \setcounter{secnumdepth}{3}


% Breche URLs um, um under-full \hbox zu verhindern
% \setcounter{biburllcpenalty}{7000}
% \setcounter{biburlucpenalty}{8000}
\usepackage[bottom]{footmisc}

% Remove padding from \fbox command
\setlength{\fboxsep}{0pt}
\setlength{\fboxrule}{1pt}

% Erstellung eines Abkürzungsverzeichnis
% Siehe auch:
% https://de.wikibooks.org/wiki/LaTeX-W%C3%B6rterbuch:_Abk%C3%BCrzungsverzeichnis
\usepackage{acronym}

% Abstand für neue Absätze
\renewcommand{\\}{\vspace*{0.5\baselineskip} \newline}

% Hurenkinder und Schusterjungen verhindern
\clubpenalty10000
\widowpenalty10000
\displaywidowpenalty=10000

% Die Abbildungen in meinem Dokument sollen unabhängig von Kapiteln und 
% Unterkapiteln durchnummeriert werden, beginnend mit 1
\renewcommand{\thefigure}{\arabic{figure}}

% Fußnoten werden durchlaufend durch das Dokument nummeriert und starten nicht bei jedem Kapitel bei 1
\usepackage{chngcntr}
\counterwithout{footnote}{chapter}
\counterwithout{figure}{chapter}

\usepackage{xcolor}
\definecolor{turkis}{HTML}{389ab5}
\definecolor{white}{HTML}{ffffff}
%\definecolor{turkis}{RGB}{45,72,80}

\usepackage{setspace}

\usepackage[font={footnotesize},labelsep=colon]{caption}

% Zur Seitenzahl, /newpage nach Roman Counting, sonst wird schon in Arabic gezählt

% Funktioniert anscheinend nicht in Kombi mit hyperref
%\usepackage{titleref}

\setlength\bibitemsep{15pt}  % Abstand zwischen 2 Einträgen im LitVZ
%\setlength{\bibhang}{2em} % Einrücken 2. Zeile im LitVZ

% \setlength\itemsep{0mm}

\setlist{noitemsep}